% !PRINTED none
% Abschnitt-title page, no running folio visible in scan: "Einundzwanzigster Abschnitt. Die Normen der Klasseninvarianten f(omega)." heading, rule, then § 140 heading and body text begin.

\begin{center}
\gesperrt{Einundzwanzigster Abschnitt.}

\textbf{Die Normen der Klasseninvarianten $f(\omega)$.}

\rule{0.3\linewidth}{0.4pt}
\end{center}

\begin{center}
§ 140. \gesperrt{Konvergenz einer unendlichen Reihe.}
\end{center}

Wir betrachten eine beliebige quadratische Form
\weq{1}{\psi(x,y) = Ax + 2Bxy + Cy^2} % [sic] source prints $Ax$; the quadratic form requires $Ax^2$
mit negativer Diskriminante
\weq{2}{-m = B^2 - AC,}
in der vorläufig die Koeffizienten $A$, $B$, $C$ nicht notwendig als ganze Zahlen betrachtet werden sollen, nur sollen sie reell sein. Nehmen wir $x$, $y$ als rechtwinkelige Koordinaten in einer Ebene an, so entsprechen den ganzzahligen Werten von $x$, $y$ die \gesperrt{Gitterpunkte}. Die Gleichung
\[\psi(x,y) = t\]
stellt für ein konstantes $t$ eine Ellipse dar, und die Anzahl $Z(t)$ der Werte von $\psi$, für \gesperrt{ganzzahlige} $x$, $y$, die kleiner als $t$ sind, ist gleich der Anzahl der im Innern dieser Ellipse gelegenen Gitterpunkte. Nach Bd.\,II, § 194, (5) nähert sich der Grenzwert des Verhältnisses $Z(t):t$ für ein unendlich großes $t$, dem Flächeninhalt der Ellipse $\psi = 1$, nämlich dem Werte $\pi:\sqrt{m}$. Also haben wir
\weq{3}{\mathrm{Lim}\, \frac{Z(t)}{t} = \frac{\pi}{\sqrt{m}}.}
Lassen wir also in der unendlichen Reihe
\weq{4}{S = \sum^{x,\,y} \frac{1}{\psi(x,y)^s}}
die Zahlen $x$, $y$ alle ganzzahligen Werte, ausgenommen die Kombination $x=0$, $y=0$, durchlaufen, so ist nach Bd.\,II, § 196, 4 $S$ für $s>1$ konvergent, und es ist
\weq{5}{\mathop{\mathrm{Lim}}\limits_{s=1}\, (s-1)S = \frac{\pi}{\sqrt{m}}.}
